% Round 17. Foundation branch. Candidate; independent review required.
% This source does not edit the odd-characteristic scope of ON17 or QF17.
\section{Nondegenerate quadratic reduction in residue characteristic two}
\label{r17c2:section}

\paragraph{Scope and classical antecedents.}
\textbf{Internally reviewed.} The final claims register identifies the independent audits and frozen dependencies.
We prove the characteristic-two version needed for odd-dimensional
ordinary quadratic reduction. The exact quadratic change of
coordinates uses an explicit etale isometry slice, rather than division
by two or extraction of square roots of units. We also prove
the quadric intersection and cycle-class inputs in every
characteristic. The resulting local and height theorems retain
native residue fields, all field degrees and all ordinary actual
tests. The even-dimensional degenerate quadratic case in residue
characteristic two is not included in this section.

Jannsen--Saito, \emph{Bertini theorems and Lefschetz pencils over
discrete valuation rings}, Section 3, Lemma 4(i) and Theorem 5(c),
\url{https://arxiv.org/pdf/0911.1470}, state the classical
nondegenerate normal form and its order reduction, including
residue characteristic two. The new use here is their activation
through a native all-degree numerical correction and the frozen
ordinary actual detector. No priority claim is made for the
normal-form or quadric facts.

\subsection{Quadric cycle comparisons without a characteristic restriction}

\begin{lemma}[Split quadric cells and native numerical comparison]
\label{r17c2:quadric-cycles}
Let $E/k$ be a smooth projective quadric over a finite field
of arbitrary characteristic, and let $\ell\ne\operatorname{char}k$.
Every native numerically zero rational cycle on $E$ has zero
geometric $\ell$-adic class. The native numerical groups are
one-dimensional ample lines except possibly in the middle
codimension of an even-dimensional quadric. In that middle
degree their dimension is at most two.
\end{lemma}

\begin{proof}
A smooth quadric splits over some finite extension $k'/k$.
For completeness, over the algebraic closure one splits off
successive hyperbolic pairs: an isotropic vector in a smooth
quadric is not in the polar radical, and a vector with polar
pairing one can be modified by a multiple of it to become
isotropic. Repeat on the orthogonal complement. In odd vector
dimension the final one-dimensional summand is a square term.
The finitely many chosen basis vectors and coefficients are
defined over a finite extension of the finite field. This
argument also applies in characteristic two to smooth odd
dimensional quadrics with their one-dimensional anisotropic
polar radical.

We record the split calculation without importing a
characteristic-not-two cellular theorem. Write $Q_d$ for a
split smooth projective quadric of dimension $d$. For $d\ge2$
choose a hyperbolic pair, so its equation is
\[
                  x_0x_1+q'(x_2,\ldots,x_{d+1})=0.
 \tag{C217.1}\label{r17c2:split-equation}
\]
The open set $x_0\ne0$ is $\mathbf A^d$. The reduced
closed complement is the projective cone over $Q_{d-2}$,
with vertex $[0:1:0:\cdots:0]$. Removing that vertex is
a line bundle over $Q_{d-2}$. A closed affine-cell filtration
of the base therefore gives one on the cone: include the
vertex first and then the cones over the successive closed
subsets. Each new stratum is a line bundle over an affine
cell and hence is affine space of dimension one larger.
Finally adjoin the open $\mathbf A^d$.
The bases are $Q_0$, two rational points, and $Q_1$, the
split smooth conic $\mathbf P^1$. In characteristic two
a tangent-hyperplane complement on that conic can be
nonreduced; the filtration uses its reduced point, which
has exactly the same open complement. Chow localization
and etale cohomology use these reduced closed supports.

This gives one affine cell in every dimension from $0$
through $d$, with a second cell in dimension $d/2$ if
$d$ is even. The ordinary Chow localization sequence
shows that the closures of the cells generate the rational
Chow groups. In etale Borel--Moore cohomology the localization
sequence and the single even-degree class of an affine cell
show inductively that those closures form a basis: there
are no odd terms which could produce boundary relations.
Compatibility of the cycle class with localization gives
\[
 \operatorname{CH}^j(Q_d)_{\mathbf Q}\otimes\mathbf Q_\ell
        \xrightarrow{\sim}
 H^{2j}((Q_d)_{\bar k},\mathbf Q_\ell(j)),
 \tag{C217.2}\label{r17c2:split-cycle-iso}
\]
and all odd cohomology is zero. Proper smooth duality makes
the corresponding cycle intersection pairing nondegenerate.

If a native cycle $z$ on $E$ is numerically zero, its
pullback to $E_{k'}$ remains numerically zero. Indeed for
every complementary $\xi$ upstairs, finite projection gives
\[
 \deg_{k'}(z_{k'}\xi)
      =\frac1{[k':k]}\deg_k\bigl(z\,\operatorname{pr}_*\xi\bigr)=0.
\]
The split cycle isomorphism and its perfect intersection
pairing make its geometric class zero. This is also the
original geometric class of $z$. Restriction to $k'$ is
injective on native numerical groups by projection and
the full nonzero finite degree. The split dimensions just
computed give the asserted dimension bounds, and the
nonzero native ample powers fill each one-dimensional
off-middle group.
\end{proof}

\begin{lemma}[The quadric middle form in every characteristic]
\label{r17c2:quadric-form}
Let $E/k$ be a smooth projective quadric of dimension
$2a\ge2$ over a finite field of arbitrary characteristic.
For every ample rational divisor $h$, the form
\[
                (-1)^a\deg_k(u^2)
\]
is positive definite on the real extension of
$\{u\in N^a(E)_{\mathbf Q}:\deg_k(u h^a)=0\}$.
\end{lemma}

\begin{proof}
Pass to a finite splitting field only for the numerical
calculation. In the hyperbolic equation
$\sum_{i=0}^a x_i y_i=0$, cut by $x_1=\cdots=x_a=0$.
The resulting cycle is the union of two maximal planes,
of classes $l,l'$, so for the natural hyperplane $h_0$,
\[
                         h_0^a=l+l'.
\]
This is an actual complete-intersection cycle identity,
also in characteristic two. For either maximal plane
$\mathbf P^a\subset E$, smoothness gives the normal-bundle
sequence
\[
 0\longrightarrow N_{\mathbf P^a/E}
   \longrightarrow\mathcal O(1)^{\oplus(a+1)}
   \longrightarrow\mathcal O(2)\longrightarrow0.
\]
Its top Chern coefficient is
\[
 \deg(l^2)=\deg((l')^2)
     =[t^a]\frac{(1+t)^{a+1}}{1+2t}
     =\frac{1+(-1)^a}{2}.
\]
Since $\deg(h_0^a l)=1$, it follows that
\[
       \deg(ll')=\frac{1-(-1)^a}{2},\qquad
       \deg((l-l')^2)=2(-1)^a.
 \tag{C217.3}\label{r17c2:middle-matrix}
\]
All coefficients are integer intersection numbers. Their
use does not require that the characteristic invert two.
The displayed matrix is nondegenerate, so these two plane
classes form the middle numerical basis supplied by the
cell calculation.

For $a\ge2$, the geometric Picard group has rank one by
the same cell calculation. Thus an ample $h$ is a positive
rational multiple of $h_0$ numerically, and its orthogonal
middle line is the line of $l-l'$. For $a=1$, the split
quadric surface is $\mathbf P^1\times\mathbf P^1$ in
every characteristic. Write $h=\alpha l+\beta l'$ with
$\alpha,\beta>0$. Its orthogonal line is generated by
$\alpha l-\beta l'$, whose square is $-2\alpha\beta$.
These calculations prove definiteness after splitting.

Native numerical restriction is injective, and
$\deg_{k'}(u_{k'}v_{k'})=\deg_k(uv)$ when each degree is
taken over its own base field. Restrict the definite form
to the native subspace. This proves the lemma without
putting an ordering on any $\mathbf Q_\ell$-space.
\end{proof}

\begin{proposition}[Quadric forests in arbitrary residue characteristic]
\label{r17c2:forest}
The numerical equality $\ker A_{n-r}=\ker\rho_{n-r}$ and
the all-degree strong-G conclusion of QF17 hold for a
supplied reduced strictly semistable native quadric forest
in every residue characteristic. All its other hypotheses,
including $n\ge3$, remain in force.
\end{proposition}

\begin{proof}
The incidence formula and the cut-edge construction of
independent ample powers in QF17 use only the reduced
principal fiber and the fact that the graph is a forest.
In nonmiddle degree the groups on every double quadric
are the one-dimensional ample lines of
Lemma~\ref{r17c2:quadric-cycles}; complementary independent
edge powers annihilate every edge difference of an
adjoint-kernel tuple.

In middle degree $n=2a+1$, cut-edge ample powers make
each difference primitive for the common ample class.
The identity
$0=-\sum_e\deg_k((\rho_a v)_e^2)$ and
Lemma~\ref{r17c2:quadric-form} force every difference to
vanish, with the positive native weights $[k_e:k]$ kept.
No component positivity is assumed. This proves the same
kernel equality.

Lemma~\ref{r17c2:quadric-cycles} turns numerical equality
of actual restriction differences into geometric
cohomological equality. The reduced compatible-kernel
criterion KC17 then produces a rational numerical
admissible extension in every degree. Its cohomological
weight, native trace and numerical-duality proof has no
odd-characteristic premise. The frozen numerical detector
and exact-support reconstruction supply the full ordinary
actual statement over the native valued field.
\end{proof}

\subsection{An exact etale Morse construction with unimodular polar form}

\begin{lemma}[A pointed smooth isometry slice]
\label{r17c2:isometry-slice}
Let $R$ be a complete DVR of mixed characteristic, with
any residue characteristic. Let $U/R$ be smooth, let
$a:\operatorname{Spec}R\to U$ be a section, and let
$q_x(v)$ be a quadratic form in $N$ indeterminates $v$,
with coefficients regular on $U$. Put $q_0=q_x|_a$.
Assume the polar bilinear matrix $B_0$ of $q_0$ is
invertible over $R$. Then there is a pointed etale
neighborhood $U'\to U$, with a lift of the section and
unchanged closed residue field, and a regular matrix
$M:U'\to\operatorname{GL}_N$ such that
\[
                  q_0(Mv)=q_x(v),\qquad M|_a=1.
 \tag{C217.4}\label{r17c2:isometry}
\]
\end{lemma}

\begin{proof}
The coefficient equations $q_0(Mv)-q_x(v)=0$ inside
$U\times\operatorname{GL}_N$ have the identity solution
over $a$. The derivative in the matrix direction $A$
at that solution is the quadratic polynomial
\[
                   B_0(v,Av)=v^{\mathsf T}B_0Av.
 \tag{C217.5}\label{r17c2:differential}
\]
Because $B_0$ is invertible, $C=B_0A$ can be any matrix.
Its resulting quadratic coefficients are $C_{ii}$ on
$v_i^2$ and $C_{ij}+C_{ji}$ on $v_iv_j$. Taking $C$
upper triangular makes these arbitrary, even in
characteristic two. This proves smoothness of the
isometry scheme at the identity and also yields a
particularly explicit etale slice, as follows.

Let $s$ be the $d=N(N+1)/2$ entries of an upper triangular
matrix $C(s)$, and set
\[
                M(s)=1+B_0^{-1}C(s).
\]
On the open set where $M(s)$ is invertible, map $s$ to
the $d$ coefficients of $q_0(M(s)v)$. At $s=0$ its
value is the coefficient vector of $q_0$, and the
derivative is exactly the identity: the derivative
polynomial is $v^{\mathsf T}C(s)v$, whose coefficients
are the independent upper triangular entries themselves.
Thus this coefficient map is etale on a neighborhood of
the entire zero section over $R$.

Base change that etale map by the coefficient map
$x\mapsto q_x$ on $U$. The resulting $U'\to U$ is
etale and has the lift of $a$ given by $s=0$, so its
selected closed point has unchanged residue field.
Its matrix $M(s)$ satisfies the exact coefficient
equations and is the identity along $a$. This proves
the assertion without square-root extraction, an
unproved smooth-section choice, or a formal series.
\end{proof}

\begin{proposition}[Native normal form when the Hessian is invertible]
\label{r17c2:normal-form}
Let $\mathcal X/R$ be flat of finite type, with smooth
generic fiber. Suppose a special-fiber ordinary quadratic
point $p$ of relative dimension $n$ has invertible polar
Hessian on its $(n+1)$-dimensional tangent space.
Let $R_p/R$ be the native unramified extension for
$\kappa(p)$. Then there is a common residue-isomorphic
pointed etale neighborhood with
\[
 \operatorname{Spec}R_p[z_0,\ldots,z_n]/(Q(z)-u\pi^b),
       \qquad u\in R_p^\times,\quad b\ge1,
 \tag{C217.6}\label{r17c2:normal-equation}
\]
where the polar matrix of $Q$ is unimodular over $R_p$.
In particular this applies to every ordinary quadratic
point in odd relative dimension in residue characteristic
two.
\end{proposition}

\begin{proof}
The Tor, local Nakayama and Jacobian argument of ON17's
smooth-ambient lemma does not use invertibility of two.
It gives a smooth ambient $W/R_p$ of relative dimension
$N=n+1$ and a single Cartier equation $f$. The native
gradient map has invertible Jacobian at $p$ by the stated
polar-Hessian hypothesis. Henselian lifting therefore gives
its critical section $a$ over $R_p$, just as in ON17.
Choose translated etale coordinates $x_i$ for its ideal.
For $c=f(a)$ one has $f-c\in(x)^2$; moreover
$0\ne c\in\pi R_p$, since otherwise the generic
hypersurface would have a critical point on it.

Write the exact expression
\[
        f-c=\sum_{i\le j}a_{ij}(x)x_ix_j=q_x(x).
\]
This requires no division by two. Its value $q_0=q_x|_a$
has the prescribed invertible polar matrix. Apply
Lemma~\ref{r17c2:isometry-slice}, and put $z=Mx$.
Then $q_0(z)=f-c$ exactly. Along $a$ one has $x=0$
and $M=1$, so $dz=dx$. Thus $z$ is an etale coordinate
map on a pointed neighborhood. Writing $-c=u\pi^b$
and restricting this map to $f=0$ gives the stated
common pointed etale neighborhood. The only constant
extension used for this local chart is the native
unramified $R_p/R$, with unchanged residue field at
the selected point.

Finally, in characteristic two with $N$ even, the polar
matrix is alternating. If it were degenerate, its radical
would have dimension at least two. Over the algebraic
closure the restriction of a quadratic form to this
radical is the square of a linear form, and hence vanishes
at a nonzero radical vector. That vector would be a
singular projective point of the tangent quadric. Its
smoothness rules this out. Thus ordinary quadratic points
with $n$ odd have exactly the invertible Hessian required.
\end{proof}

\subsection{The native semistable chains in residue characteristic two}

\begin{proposition}[Order reduction for an invertible polar quadratic form]
\label{r17c2:chains}
Assume $R=K^\circ$ for $K/\mathbf Q_2$ and $n\ge3$ is
odd. Let an integral flat projective model with smooth
geometrically integral generic fiber have geometrically
integral special fiber, smooth except at finitely many
ordinary quadratic points. Make the single extension
$K'=K(\tau)$ with $\tau^2=\pi$.
At a native point with normal-form order $b$, exactly
$b$ successive blowups at the remaining reduced point
give the chain of QF17: $b-1$ projective cone resolutions
$\mathbf P_E(\mathcal O_E\oplus\mathcal O_E(-1))$ and
one terminal smooth quadric. The resulting global model
is regular, projective and reduced strictly semistable,
with a native quadric tree and no triple strata.
\end{proposition}

\begin{proof}
Proposition~\ref{r17c2:normal-form} supplies exact native
equations $Q(x)=u\pi^b$, with invertible polar matrix.
The extension $\tau^2=\pi$ is Eisenstein of full degree
two and residue degree one, even though it is wild.
The equation becomes $Q(x)=u\tau^{2b}$.

The saturated reduced-point blowup charts are exactly
\[
 x_i v=\tau,\qquad
 Q(y_0,\ldots,1,\ldots,y_n)
                          =u x_i^{2b-2}v^{2b},
 \tag{C217.7}\label{r17c2:x-chart}
\]
and
\[
                         Q(z)=u\tau^{2b-2}.
 \tag{C217.8}\label{r17c2:tau-chart}
\]
The Rees saturation verification in QF17 uses only a
nonzero polynomial modulo the chart generator and
permutation of a regular sequence in the semistable
Cohen--Macaulay ambient chart. It uses neither two as a
unit nor tame ramification, so it applies unchanged.
The $\tau$-chart polynomial is primitive over its DVR,
giving the same injectivity of multiplication by $\tau$.

We check every smoothness step which previously invoked
odd characteristic. Write $\overline B$ for the invertible
polar matrix over the native residue field. For a vector
$X$ with $X_i=1$, its gradient is $\overline B X$.
Euler's polynomial identity still gives
$\sum_jX_j\partial_j\overline Q=2\overline Q=0$
in characteristic two. If every derivative in a $y_j$
direction vanished, it would force the remaining $i$
derivative to vanish as well. This contradicts
$\overline B X\ne0$. Thus one $y$ derivative is a unit
at every closed-fiber point in \eqref{r17c2:x-chart},
also when $b=1$ and the right side is $\overline u v^2$.
Eliminating that variable etale locally leaves precisely
$x_i v=\tau$. These charts are regular strictly
semistable with multiplicity one.

For $b>1$, the only remaining singular special point of
\eqref{r17c2:tau-chart} is $z=0$, since its gradient
$\overline Bz$ vanishes only there. Its order has decreased
by two. For $b=1$, the equation is $Q(z)=u$ and is smooth
over the DVR: $z=0$ is not on its special fiber, so its
gradient cannot vanish. The terminal exceptional quadric is
\[
                  T=\{\overline Q(X)=\overline u S^2\}
                                      \subset\mathbf P^{n+1}.
 \tag{C217.9}\label{r17c2:terminal}
\]
It is smooth in characteristic two. Indeed all its
$X$ derivatives vanishing would imply $X=0$, and its
equation would then imply $S=0$, impossible projectively.
Its section $S=0$ is the smooth double quadric
$E=\{\overline Q=0\}\subset\mathbf P^n$.

For larger order the exceptional component is the cone
over this $E$. Its vertex blowup is the same projective
bundle as in QF17, with its two disjoint sections. This
identification follows from the line-to-vertex map and
its direction charts, without a characteristic hypothesis.
The number of blowups and components follows by induction.
The native pointed etale charts identify the associated
graded rings, so each exceptional variety is defined over
the original point's residue field, not a chosen geometric
splitting field. Blowup commutes with the etale maps.

Away from the original finitely many points the model is
smooth, and their chains are disjoint. Consequently the
global model has the stated tree, smooth components and
double strata, and no triple points. Projectivity is
preserved by the blowups; the displayed charts and the
integral generic fiber give integrality and flatness.
This is also the nondegenerate case of the original
Jannsen--Saito order-reduction theorem.
\end{proof}

\begin{theorem}[All-degree actual G for odd-dimensional ordinary quadratic reduction at two]
\label{r17c2:actual}
Let $K/\mathbf Q_2$ and its absolute value be fixed, and
let $X/K$ be smooth projective geometrically integral of
odd dimension $n\ge3$. Suppose it has an integral flat
projective model with geometrically integral special fiber,
smooth except at finitely many ordinary quadratic points.
Then every geometrically homologically trivial rational
cycle in every positive codimension has strong ordinary
actual G on its exact specified support complement, for
every original compact actual test. No regularity of the
initial total model or bound on the node orders is required.

The conclusion holds anew for all such cycles after every
finite extension of $K$, and for every smooth projective
target admitting such a projective dominant witness after
a finite auxiliary local extension.
\end{theorem}

\begin{proof}
The chain proposition gives a regular reduced quadric
forest over $K'$. Apply the all-characteristic forest
criterion of Proposition~\ref{r17c2:forest}. It constructs
native rational numerical admissible lifts in every degree.
The frozen original-field numerical witness and dominant
Bott--Chern detector descend obstruction zero, retaining
the full auxiliary degree two and every node residue
degree. The frozen curved pair for the final specified
cycle then supplies its exact-support all-actual form.
No alteration is claimed finite on that complement.

After any further finite extension, flatness and the
integral generic fiber preserve integrality of the model.
The special geometric nodes and their invertible polar
forms persist, while each order is multiplied by the
ramification index. Thus the same theorem applies to all
newly defined cycles, not only pullbacks. The dominant
witness assertion is the frozen numerical detector, with
the full composite function-field degrees.
\end{proof}

\begin{corollary}[The ordinary quadratic height sector at all finite primes in odd dimension]
\label{r17c2:height}
Let $X/F$ be smooth projective geometrically integral over
a number field. Suppose every finite completion has a
projective dominant witness after a finite auxiliary local
extension with good or ordinary quadratic reduction. For
the ordinary quadratic alternative above residue two,
assume the witness has odd dimension at least three;
above other primes use ON17's range.
Then the canonical height is defined on the entire
geometric homological rational Chow groups in every pair
of complementary degrees. In particular this applies
when an odd-dimensional $X$ itself has good or ordinary
quadratic reduction at every finite place.
The value has the frozen number-field normalization and
the automatic arithmetic-intersection comparison of GA17.
\end{corollary}

\begin{proof}
ON17, the present theorem, and the frozen good-reduction
theorem supply local obstruction zero for every geometric
homological class at every finite place. The frozen
canonical height theorem and GA17 therefore apply to the
complete homological groups for the stated reduction
sector. No assertion is made that arbitrary varieties
have witnesses of this type.
\end{proof}

\paragraph{Remaining characteristic-two case.}
For even relative dimension in residue characteristic two,
an ordinary tangent quadric has a one-dimensional polar
radical, so the gradient-section argument and the isometry
slice above do not apply. Jannsen--Saito treat that case
by a different polynomial $P(x)+z^2+bz+c$ and a different
order reduction. Its integration into the present native
numerical proof requires a separate statement, field-degree
account and review; it is not included by analogy here.
